\documentclass[10pt,a4paper]{amsart}
\usepackage[margin=19mm]{geometry}
\usepackage{amsmath,amssymb,mathtools}
\usepackage[scr=rsfs,cal=euler]{mathalfa}
\usepackage{xfrac}
\usepackage[unicode,hidelinks,bookmarks=false]{hyperref}
\newcommand{\ie}{i.e.\ }
\newcommand{\cf}{cf.\ }
\newcommand{\resp}{respectively\ }
\newcommand{\Subsection}{\subsection}
\providecommand{\nobreakdash}{\nobreak\hbox{-}}
\setlength{\emergencystretch}{2em}
\multlinegap0pt
\usepackage{bm}
\usepackage{amssymb}
\usepackage{tikz}
\usepackage{enumerate}
\usepackage[shortlabels]{enumitem}
\usepackage{float}
\newtheorem{claim}{Claim}
\usepackage{cleveref}
\crefname{claim}{Claim}{Claims}
\newcommand{\R}{{\mathbb {R}}}
\newcommand{\cosec}{\mathrm{cosec}}
\title{On a variant of bilinear spherical maximal function}
\author{Ankit Bhojak, Surjeet Singh Choudhary, Saurabh Shrivastava}
\date{Source: arXiv:2410.17583v2 (CC BY 4.0)}
\begin{document}
\maketitle

\noindent\emph{Source and licence.} On a variant of bilinear spherical maximal function. Version arXiv:2410.17583v2 (CC BY 4.0), distributed by arXiv under the Creative Commons Attribution 4.0 International licence, http://creativecommons.org/licenses/by/4.0/, as recorded in the arXiv licence metadata for this exact version and shown on the abstract page https://arxiv.org/abs/2410.17583v2. Full licence notice: \texttt{licenses/arxiv-2410.17583-CC-BY-4.0.txt}.\par\smallskip

\noindent\textit{Editorial scope.} Counted unit: the article's claim claim:sublevel (source0.tex lines 1186--1191). The article's complete original proof of it is delivered with it (source0.tex lines 1192--1220, one \texttt{proof} environment of 29 lines). No mathematics is written, repaired or re-derived: the statement and the proof are the article's own lines, in the article's own notation, and no retained line is substituted or re-flowed.  No further source-local context is needed: every cross-reference of every retained line resolves inside the two delivered spans.  Nothing was omitted from any retained span, so the package carries no omission record.  No retained line contains a citation, so the wrapper carries no bibliography and no bibliography entry.  No retained line contains a figure environment, an image-inclusion command or drawing code, so no image is delivered and the package declares no asset dependency.  Editorial formatting only, in addition: an AMS \texttt{amsart} preamble that restates the article's own declarations and macros used by the retained lines, namely the environments \texttt{claim} on separate counters, together with the standard packages those lines need.  The retained environments print in this document as Claim 1 in order of appearance, whereas the article prints the same items with the numbering of its own full document.  The complete native source of the article as distributed in the arXiv e-print for this exact version is bundled unchanged as \texttt{sources/167284/source0.tex}.
	\begin{claim}\label{claim:sublevel}
		For $(\bar u,\bar v)\in\Sigma_B$, we have
		\begin{equation*}
			|\{w\in I_B^{(\bar u,\bar v)}:\;|J(\boldsymbol{T}^{-1}(\bar u,\bar v,w))|\leq C\delta\epsilon\}|\lesssim\delta^3\epsilon^{1-2a}.
		\end{equation*}
	\end{claim}

	\begin{proof}[Proof of \Cref{claim:sublevel}:]
		Let $G(w)=J(\boldsymbol{T}^{-1}(\bar u,\bar v,w))=\alpha({\bar z})\cot(t_1)\tan(t_2)\cot(t_3)-\beta({\bar z}+(\bar u,\bar v))$. We observe that the following holds by the definition of $\boldsymbol{T}$.
		\[\frac{\partial t_1}{\partial w}=1,\quad\frac{\partial t_2}{\partial w}=\frac{\sin(t_1-t_3)}{\sin(t_2-t_3)},\quad\text{and}\quad\frac{\partial t_3}{\partial w}=\frac{\sin(t_1-t_2)}{\sin(t_2-t_3)}.\]
		Thus, we have that
		\begin{align*}
			|G'(w)|=&\bigg|\alpha({\bar z})\Big(-\cosec^2(t_1)\tan(t_2)\cot (t_3)\frac{\partial t_1}{\partial w}+\cot(t_1)\sec^2(t_2)\cot(t_3)\frac{\partial t_2}{\partial w}\\
			&\hspace{1mm}-\cot(t_1)\tan(t_2)\cosec^2(t_3)\frac{\partial t_3}{\partial w}\Big)\bigg|\\
			=&\bigg|\frac{\alpha({\bar z})}{4\sin^2(t_1)\cos^2(t_2)\sin^2(t_3)}\widetilde{J}(t_1,t_2,t_3)\bigg|,
		\end{align*}
		where the function $\widetilde{J}:I^3\to\R$ is defined as
		\begin{align*}
			\widetilde{J}(t_1,t_2,t_3)=&\frac{1}{\sin(t_2-t_3)}\big[\sin (2t_2)\sin(2t_3)\sin(t_2-t_3)+\sin (2t_1)\sin(2t_2)\sin(t_1-t_2)\\
			&\hspace{3cm}+\sin (2t_1)\sin(2t_3)\sin(t_3-t_1)\big].
		\end{align*}
		We have the following identity,
		\begin{align*}
			&\big[\sin (2t_2)\sin(2t_3)\sin(t_2-t_3)+\sin (2t_1)\sin(2t_2)\sin(t_1-t_2)\\
			&\hspace{2cm}+\sin (2t_1)\sin(2t_3)\sin(t_3-t_1)\big]\\
			=&-2 \sin \left(\frac{t_1-t_2}{2}\right) \sin \left(\frac{t_2-t_3}{2}\right) \sin \left(\frac{t_3-t_1}{2}\right)\\
			&\hspace{2cm}\times\big[\cos (t_1+t_2-2 t_3)+\cos (t_1-2 t_2+t_3)+\cos (-2 t_1+t_2+t_3)\\
			&\hspace{2.5cm}+\cos (2 t_1+t_2+t_3)+
			\cos (t_1+2 t_2+t_3)+\cos (t_1+t_2+2 t_3)\\
			&\hspace{2.5cm}+\cos (2 (t_1- t_2))+\cos (2 (t_2-t_3))+\cos (2 (t_3-t_1))\\
			&\hspace{2.5cm}+2 \cos (t_1-t_2)+2 \cos (t_2-t_3)+2 \cos (t_3-t_1)+1\big].
		\end{align*}
		Using the above identity it is clear that $|\widetilde{J}(t_1,t_2,t_3)|\gtrsim|t_1-t_2||t_3-t_1|\gtrsim\delta^2\epsilon^{2a}$. Using this and $|\alpha({\bar z})|\sim1$ (this is the only place where we require this hypothesis), we obtain that
		\[|G'(w)|\gtrsim\delta^{-2}\epsilon^{2a}.\]
		Hence, the above estimate and an application of mean value theorem produces the required claim.
	\end{proof}

\end{document}
